\chapter{DAS系统验证与结果分析}

\section{引言}

第二章建立了外差式$\varphi$-OTDR的瑞利散射、距离映射、平衡相干探测和数字相位解调模型，第三章完成了无EDFA分光集成器件发射链路、SOA--AOM--DAQ协同门控以及DAS一体化收发模组的光电功能验证。在此基础上，本章进一步从端到端传感角度评价系统对已知振动的恢复能力。验证对象不再是单个器件或中间节点，而是由探测脉冲发射、光纤瑞利散射、双偏振外差接收、同步采集和数字相位解调共同构成的完整测量链路。

本章沿用第三章确定的统一硬件工作点：LD电流和TEC温度分别为200~mA和30~$^\circ$C，SOA脉冲电流为1600~mA，脉宽和重复频率分别为100~ns和2~kHz，关断偏压为$-3$~V；AOM采用80~MHz移频，并以80~ns提前量和20~ns保持量实施预补偿门控。上述参数在振动实验中保持固定，使后续结果对应同一发射功率、脉冲宽度、外差频带和采集时序，避免将硬件工作点变化混入解调及系统比较结果。

接收端同步记录P、S两路外差拍频。数据处理采用第二章所述统一流程：首先校正两通道的直流偏置、幅值比例、I/Q方向和相对时延，再以固定数字本振完成正交下变频和低通滤波，构造复包络并进行复共轭空间差分；随后依据幅值可信度实施双偏振功率加权合成，获得沿光纤距离和慢时间分布的差分相位。通过对已知位置、频率和驱动波形的压电陶瓷振动进行测量，从扰动定位误差、频率误差、谱峰信噪比、波形相关性和有效数据率等指标评价系统性能。所有数据采用相同的距离映射、滤波、差分标距和无效点判据处理。

本章首先给出传感光纤、压电陶瓷扰动、采样触发及数据规模等实验条件；随后完成双偏振通道标定，并依次报告扰动位置、频率和波形恢复结果。在此基础上，以第三章定义的分立系统D和集成系统M为比较对象，在相同传感光纤、PZT输入、采集设置和数字处理流程下开展交替测试，评价模组集成对端到端传感结果的影响。最后通过独立上电重复测量和连续运行记录，分析定位、频率、波形相关性及有效数据率在实验时长内的重复性和漂移。由此形成从模组光电输出到振动相位恢复的系统级验证，并为第五章总结本文收发模组的适用条件和工程特性提供依据。

\section{实验条件与数据采集}

系统总体结构与器件实现已在第三章给出，本章沿用第三章建立的SOA--AOM--DAQ统一时序。本节给出振动验证实验的具体条件，包括测试光纤与PZT扰动设置，以及采样率、触发与数据规模。

\begin{writingplan}
  \item 给出传感光纤长度、PZT位置与驱动参数，以及D/M测试共用的扰动和环境条件。
  \item 列出采样率、记录长度、触发和SOA--AOM--DAQ时序，并给出距离映射、测距范围与空间分辨率。
  \item 用一幅实验布置图和一张参数表汇总上述条件及数据处理范围。
\end{writingplan}

\section{信号解调与振动恢复}
\label{sec:chap4_channel_calibration}
\label{sec:chap4_fixed_demod}

数字下变频、I/Q提取、复共轭差分与双偏振合成的原理已在第二章给出，本章按该统一流程处理采集的数据。本节在给出通道标定与解调参数配置后，报告扰动定位、频率和波形恢复结果。

\begin{writingplan}
  \item 用P、S两路原始波形及暗态/断光结果确定偏置、增益、I/Q方向和相对时延校准量，并检查ADC饱和与校准残差；本振频率、滤波带宽、差分标距、合成规则及无效点处理以流程图和参数表集中给出，D/M数据使用相同参数和代码路径。
  \item 由距离--时间相位图确定扰动位置，以无扰动区为背景，报告定位误差、空间展宽和扰动信噪比。
  \item 从定位区域提取相位序列，按统一FFT规则报告恢复频率、误差、谱峰信噪比和谐波。
  \item 将恢复波形与PZT参考波形对齐，报告相关系数、时延和连续性；用时空图、频谱及波形对比图集中展示结果。
\end{writingplan}

\section{分立与集成系统传感性能对比}

第三章已经给出分立式与集成式系统的器件边界、接口方式及模组级光电验收条件，本节在相同传感光纤、采集设置和统一解调流程下比较二者的端到端振动传感结果。

\begin{writingplan}
  \item 在同一光纤、PZT输入、采集参数和解调流程下交替测试D与M，配对比较定位误差、谱峰信噪比、频率误差、波形相关性和有效数据率。
  \item 用同尺度图和汇总表展示均值、标准差及差值不确定度，并结合第三章的发射工作点、协同门控和模组接口条件解释差异。
  \item 结论仅判断模组是否保持端到端测量能力及其收益与代价，不预设集成系统全面优于分立系统。
\end{writingplan}

\section{测量重复性与长期稳定性分析}

本节在第三章模组级功能验收基础上，评价振动测量结果本身的重复性及实验时长内的漂移。

\begin{writingplan}
  \item 采用统一预热时间和D/M交替顺序完成至少3轮独立采集，统计定位、频率、波形相关性和有效数据率的重复性。
  \item 在固定PZT输入下连续记录主要传感指标、接收光功率与温度，报告实验时长内的漂移和突变，不外推现场寿命。
  \item 用重复测量误差条和随时间曲线汇总结果，异常数据按预设规则处理并保留记录。
\end{writingplan}

\section{本章小结}

\begin{writingplan}
  \item 汇总定位误差、频率误差、波形相关性、D/M差异及实验时长内稳定性等核心结果。
  \item 说明数字处理作为统一的验证流程，并给出数据支持的系统能力与局限。
\end{writingplan}
